\section{The Hodge-atom specialization}
\label{sec:hodge}

We now apply the preceding construction to the standard Hodge atoms of
Katzarkov--Kontsevich--Pantev--Yu.  Only the data needed for the type map are
recalled; the non-archimedean A-model comparison theorems remain geometric
inputs.

For a smooth complex projective variety $X$, let $B_X^{\mathrm{Hod}}$ be
the Hodge-fixed locus in the base of its maximal non-archimedean A-model
$F$-bundle.  On the open set $U_X\subset B_X^{\mathrm{Hod}}$ where the
reduced Euler spectral cover is unramified, write
\[
 \widetilde U_X\longrightarrow U_X
 \tag{3.1}\label{eq:hodge-cover}
\]
for that finite cover.  Its connected components are the local Hodge atoms,
and the multiplicity of a component is its degree over $U_X$
\cite[Definition~5.10 and Section~5.4]{KKPY}.  Accordingly, set
\[
 A_X:=\pi_0(\widetilde U_X),
 \qquad
 m_X(a):=\deg(\widetilde U_{X,a}/U_X).
 \tag{3.2}\label{eq:hodge-occurrences}
\]
Automorphisms of $X$ act on this weighted set.  In the construction of
Section~\ref{sec:ledger}, their orbits are imposed by isomorphism
correspondences rather than silently discarded.

The required provider record is as follows.
\begin{center}
\begin{tabular}{@{}>{\raggedright\arraybackslash}p{0.22\textwidth}
                    >{\raggedright\arraybackslash}p{0.36\textwidth}
                    >{\raggedright\arraybackslash}p{0.30\textwidth}@{}}
\toprule
Operation & $F$-bundle decomposition & Occurrence consequence \\
\midrule
Disjoint union & external direct sum & disjoint union of covers \\
Blowup in codimension $c$ & ambient summand plus $c-1$ center summands &
  \eqref{eq:blowup} \\
Rank-$r$ projective bundle & $r$ base summands &
  \eqref{eq:projective} \\
\bottomrule
\end{tabular}
\end{center}
The spectral decomposition is \cite[Theorem~4.1]{KKPY}.  The blowup row is
\cite[Theorem~4.5]{KKPY}, based on Iritani's blowup decomposition
\cite[Theorem~5.18]{Iritani}.  The projective-bundle row is
\cite[Theorem~4.11]{KKPY}, based on the Iritani--Koto decomposition
\cite[Theorem~5.1]{IritaniKoto}.  These theorems identify the relevant
spectral covers over connected nonempty domains and preserve the degrees of
their components.  Expanding equal degrees supplies occurrence bijections;
different choices differ only by the internal relabellings already imposed in
Section~\ref{sec:ledger}.

\begin{theorem}[Hodge atoms are a ledger specialization]
\label{thm:hodge}
For the provider \eqref{eq:hodge-occurrences}, the set
\(H=\pi_0(\Pi)\) of Section~\ref{sec:ledger} is the abstract Hodge-atom set
\(\HAtoms\) of \cite[Definition~5.16 and Section~5.4]{KKPY}.  Under this
identification,
\[
 \mathcal B(X)=\sum_{\alpha\in\HAtoms}m_\alpha(X)[\alpha]
 =:\CF_{\mathrm{Hdg}}(X)
 \tag{3.3}\label{eq:chemical-formula}
\]
is its Hodge-atom chemical formula, and
\(H_{\le e}=\HAtoms_{\dim\le e}\).
\end{theorem}

\begin{proof}
The objects of \(\Pi\) are precisely the local components of
\eqref{eq:hodge-cover}, expanded by their degrees.  After the internal
occurrence labels are identified, its remaining generating arrows are the
isomorphism, disjoint-union, blowup, and projective-bundle elementary
correspondences used to define abstract atoms in
\cite[Sections~5.2.3--5.2.6]{KKPY}.  Passing to connected components of the
thin groupoid is therefore the same generated quotient.  Summing its
occurrences gives \eqref{eq:chemical-formula}.  Finally, both dimension
filtrations ask whether the same abstract class has a representative on a
smooth projective variety of dimension at most $e$.
\end{proof}

Combining Theorems~\ref{thm:hodge} and
\ref{thm:birational-quotient} recovers the Hodge-atom non-rationality
criterion \cite[Proposition~5.30]{KKPY}.  The ledger proof additionally
exhibits the universal fold: every commutative-monoid-valued statistic of
Hodge-atom occurrences that is constant on the elementary correspondences
factors uniquely through \(\mathcal L\).

\begin{remark}[Abstract and geometric atoms]
\label{rem:two-quotients}
The geometric atomic $F$-bundles of \cite[Definition~5.21]{KKPY} form a
different quotient.  Their Proposition~5.22 gives a map from abstract atoms
to geometric atomic $F$-bundle classes, but does not assert that it is
injective.  The groupoid \(\Pi\), the chemical formula
\eqref{eq:chemical-formula}, and the dimension filtration use abstract
classes.  Consequently an invariant first defined on geometric $F$-bundle
classes becomes an atom marker only after its invariance under every
elementary correspondence has been proved.
\end{remark}
